\section{Evaluation and Reflection}\label{sec:evaluation}

\subsection{Evaluation strategy}\label{sec:evalstrategy}

The evaluation is organised around the questions of the introduction rather than
a single accuracy number, because the object of study is the certificate and its
trustworthiness. Correctness is checked before any result is interpreted: each
audited condition has a unit test, and every certificate is recomputed from its
saved ingredients. Every cell runs five seeds and is reported as a mean with a
standard deviation, the sole exception being the CIFAR-10 15-layer cell, where
two runs did not train. The two claims easiest to confound, the effect of the
prior and of difficulty, rest on controlled comparisons instead: a paired
learnt-versus-random prior on every dataset, and a label-noise sweep with
architecture and sample size held fixed. Causal claims are restricted to those.

\subsection{Threats to validity}\label{sec:threats}

The main threats concern differences from the reference protocol and limits on
how the results can be interpreted. On the first, the reduced-compute protocol of
Section~\ref{sec:protocol} deviates in several ways and only one of them, the
Monte-Carlo budget, has been measured (Section~\ref{sec:mc}); the absent
hyper-parameter search and the epoch budget can move the certificate either way
and we do not bound them. One threat is internal to our own correction: the union
over prior scales needs a candidate set fixed independently of the bound set, and
ours was written down afterwards (Section~\ref{sec:confidence}). The exposure is
small, since the grid covers every scale the probe could plausibly have returned,
but that is an assumption about our procedure, not something the data confirms.

The interpretive limits concern the certified predictor, the cross-dataset
comparison, the prior comparison and seed variability. Only the Gibbs classifier
is certified (Section~\ref{sec:scope}); a formal bound for the deterministic
predictors via the C-bound is left to future work. As noted in
Section~\ref{sec:main}, the cross-dataset comparison is descriptive because its
confounders are not controlled, and the label-noise study is where a difficulty
claim is made instead. Even the prior comparison, our most controlled effect, compares
two whole protocols rather than one isolated factor, which is why a CIFAR-10
random-prior control is included rather than inferred from the easier datasets.
Finally, the across-seed spread is the variability of the training procedure and
must not be read as a confidence interval on the certificate: the PAC-Bayes
guarantee is about one trained model, holding with probability at least
$1-\delta_{\mathrm{MC}}-\delta_{\mathrm{PAC}}=0.99$.

\subsection{Limitations}\label{sec:limitations}

A non-vacuous certificate is not necessarily a useful one: a CIFAR-10
certificate below $1$ may still be too loose to drive a deployment decision.
Depth is varied only in the CIFAR-10 $f_{\text{quad}}$ sweep and never with
width; we test no modality beyond image benchmarks, and implement neither the
differentially-private-prior nor the majority-vote route. The audit corrects and
tests the pipeline issues we identified, but is not an exhaustive verification of
the upstream code.

\subsection{Reflection on the project process}\label{sec:reflection}

The most useful lesson from the project was how easily a reproduction of a
certificate can be wrong in a way that nothing flags. The original difficulty
ladder appeared to hold, the small-data certificate appeared to stay tight, and
the objective ordering appeared to match the literature, until each was traced
back to a specific step in the pipeline. The sample-count convention is the
clearest example: at $10\%$ data, putting the full $60{,}000$ in the denominator
instead of the true $n_{\mathrm b}=3{,}000$ makes a certificate that is really
$0.152$ read as $0.092$, and either figure looks plausible. That
experience shifted the project towards verifiable certificates, and motivated
building the unit tests and the registry before the headline runs.

With hindsight, one thing would have been done much earlier: profiling. The data
pipeline, not the network, dominated wall-clock time, and the Monte-Carlo
evaluation ran in far smaller batches than the hardware needed. Fixing both left
the reported numbers untouched, since neither changes what is computed, but it
turned the full matrix from a budget question into a routine one. The registry,
the per-run provenance and the regeneration of every results figure from raw results are
what make these certificates independently checkable, and on this evidence they
belong in the result rather than after it.
